\documentclass[11pt,a4paper]{article}
\usepackage[utf8]{inputenc}
\usepackage[T1]{fontenc}
\usepackage{lmodern}
\usepackage{amsmath,amssymb,amsthm}
\usepackage{mathtools}
\usepackage{booktabs}
\usepackage{array}
\usepackage{microtype}
\usepackage{geometry}
\usepackage[ngerman]{babel}
\usepackage{hyperref}

\geometry{a4paper,margin=2.5cm,top=2.4cm,bottom=2.4cm}
\hypersetup{colorlinks=true,linkcolor=black,citecolor=black,urlcolor=black}

\newtheorem{theorem}{Satz}[section]
\newtheorem{lemma}[theorem]{Lemma}
\newtheorem{proposition}[theorem]{Proposition}
\newtheorem{corollary}[theorem]{Korollar}
\theoremstyle{definition}
\newtheorem{definition}[theorem]{Definition}
\theoremstyle{remark}
\newtheorem{remark}[theorem]{Bemerkung}
\numberwithin{equation}{section}

\title{\textbf{\texorpdfstring{Die Geometrisierung der Naturkonstanten\\
und die Grenzen der Kontinuumsbeschreibung}{Die Geometrisierung der Naturkonstanten}}\\[0.4em]
\large Axiomatische Grundlegung des EABC-Frameworks:\\
Hoffbauer-Klammer, $E_8$-Gitter und Hecke-$T_3$}
\author{\textbf{T. Hoffbauer}\thanks{Unabhängige theoretische Forschung, Bamberg, Deutschland.
E-Mail: \texttt{t.hoffbauer@posteo.de}}}
\date{14.\ September 2026}

\begin{document}
\maketitle

\begin{abstract}
\noindent
Seit Newton und Leibniz ist das stetige Raum-Zeit-Kontinuum das Arbeitsparadigma der theoretischen Physik.
Die Häufung offener Fundamentalprobleme -- die Herkunft der Feinstrukturkonstante, Divergenzen der QFT, Singularitäten der ART, die Clay-Probleme der Hydrodynamik und Eichentheorie -- markiert Grenzen dieser Beschreibung, keine bereits geschlossenen Beweise ihrer Ablösung.

Das EABC-Framework bettet die Raumzeit in die Einheits-Normschale $S^7\subset\mathbb{R}^8\cong\mathbb{O}$ des Gitters $E_8$ ein.
Der gesicherte Kern dieser Note ist die parameterfreie Form
\[
\alpha^{-1}_{\mathrm{EABC}}
=137+\sqrt{D_r^2-D_t^2},\qquad
D_r=\frac{8}{137\varphi}=\frac{4(\sqrt{5}-1)}{137},\qquad
D_t=\frac{126}{49457},
\]
mit $\varphi=(1+\sqrt{5})/2$ und $126=\tau(3)/2$.
Numerisch $\alpha^{-1}_{\mathrm{EABC}}\approx 137{,}03599954005601$.
Das Lean-4-Theorem \texttt{alpha\_inv\_hecke\_sub\_ppb} schließt das Sub-PPB-Fenster
$137{,}03599954005$ bis $137{,}03599954007$ über $\mathbb{Q}$ ohne Fließkomma-Axiome.
Der Abstand zum CODATA-2022-Zentralwert $137{,}035999177(21)$ beträgt $3{,}63\times 10^{-7}$
(etwa $17\sigma$ jenes Bandes) und darf nicht mit dem $0{,}41\,\mathrm{ppt}$-Residuum
gegen die Fenstermitte $137{,}03599954$ verwechselt werden.

Im geometrischen Einheitensystem $e=c=1$ und dem Postulat $G=4=\dim_{\mathbb{R}}\mathbb{H}$ folgen $\hbar=\alpha^{-1}$,
$\ell_P=2\sqrt{\alpha^{-1}}$, $m_P=\tfrac12\sqrt{\alpha^{-1}}$ und $h=2\pi\hbar\approx 861{,}02257886$.
Die Kapitel zu Navier--Stokes, Yang--Mills, $\Lambda$CDM, Hubble-Spannung, $g-2$ und den 25 SM-Parametern sind geometrische Lesarten und Probeformeln.
Sie sind keine Clay-vollständigen Existenzbeweise und keine Elimination der freien Parameter im Sinne einer abgeschlossenen QFT.
\end{abstract}

\noindent\textbf{MSC (2020):}
11F11, 17B25, 35Q30, 81T13, 81V10, 83C45, 03B35.\\
\textbf{Schlüsselwörter:}
EABC, Hoffbauer-Klammer, $E_8$, Feinstrukturkonstante, Hecke-$T_3$, Lean~4, Intervallarithmetik.

\tableofcontents

\section{Historische und erkenntnistheoretische Grundlagen}

\subsection{Die Illusion des Kontinuums}
Die mathematische Modellierung der physikalischen Natur ruht seit mehr als drei Jahrhunderten auf dem Dogma des Kontinuums.
Als Newton in den \emph{Principia} (1687) und Leibniz zeitgleich den Differentialkalkül formulierten, begründeten sie ein Fundament, dessen Siegeszug beispiellos war~\cite{newton1687}.
Raum und Zeit wurden als differenzierbare Mannigfaltigkeiten postuliert, auf denen infinitesimale Schritte $dt$ und $dx$ bis in jede Kleinheit getrieben werden durften.

Dieses Postulat ermöglichte klassische Mechanik, Maxwellsche Elektrodynamik, Einsteinsche Gravitation und relativistische Quantenfeldtheorie.
Doch der methodische Triumph trug den Keim einer epistemologischen Verschiebung:
der Grenzübergang $\Delta x\to 0$ war ursprünglich Rechenerleichterung, ein Kalkül zur Annäherung diskreter Summen.
Im 19.\ und 20.\ Jahrhundert wurde aus dieser Idealisierung oft ein ontologischer Anspruch:
das Kontinuum nicht mehr als Hilfsmittel, sondern als physikalische Realität selbst.

Das EABC-Framework bestreitet nicht den empirischen Erfolg der Kontinuumsphysik oberhalb atomarer Skalen.
Es bestreitet, dass der Erfolg einen Beweis für unbegrenzte Teilbarkeit von Raum, Zeit und Feld liefert.
Die vorliegende Monografie trennt deshalb scharf: was sich über $\mathbb{Q}$ einschließen lässt, und was geometrische Lesart bleibt.

\subsection{Drei Krisen der kontinuierlichen Physik}
\begin{enumerate}
\item \textbf{Ultraviolett-Katastrophe der Hohlraumstrahlung (1900).}
Im Rahmen der klassischen statistischen Mechanik verlangt das Rayleigh--Jeans-Gesetz, dass ein Hohlraumstrahler im thermischen Gleichgewicht unendliche Energiedichten im Ultravioletten abstrahle.
Die Annahme, jedem kontinuierlichen Schwingungsmodus komme dieselbe Energie $k_B T$ zu, führt zur Divergenz
\begin{equation}
U=\int_0^\infty \frac{8\pi\nu^2}{c^3}k_B T\,d\nu\longrightarrow\infty.
\end{equation}
Planck überwand dies im Dezember 1900 durch einen Bruch mit der Kontinuität: Energie nur in Quanten $E=h\nu$~\cite{planck1900}.
Damit trat die erste fundamentale diskrete Schnittgrenze in die Physik ein.

\item \textbf{Raum-Zeit-Singularitäten der Gravitation (1965--1970).}
Einstein hielt die Singularitäten der Feldgleichungen von 1915 lange für Artefakte unphysikalischer Symmetrieannahmen.
Penrose (1965) und Hawking--Penrose (1970) zeigten mit differentialtopologischen Methoden, dass unendliche Krümmung $R_{\mu\nu\rho\sigma}R^{\mu\nu\rho\sigma}\to\infty$ unter realistischen Energiebedingungen generisch ist~\cite{penrose1965,hawking1970}.
Die Singularität markiert den Zusammenbruch der Riemannschen Kontinuumsgeometrie.

\item \textbf{Divergenzen der Quantenfeldtheorie und der Landau-Pol.}
Die relativistische QFT stieß auf unendliche Selbstenergien, Vakuumpolarisationen und den Landau-Pol.
Regularisierung und Renormierung -- Cutoffs $\Lambda\to\infty$, Subtraktion von Unendlichkeiten -- lieferten endliche Messwerte.
Dirac betonte, dass dies kein Zeichen einer vollständigen Theorie sei, sondern ein Indiz, die kontinuierliche Raumzeit bei ultrakurzen Distanzen durch eine andere Kategorie zu ersetzen~\cite{dirac1975}.
\end{enumerate}

\subsection{Sommerfelds kinematische Intuition und Paulis Rätsel}
Im Jahre 1916 führte Sommerfeld zur Feinstruktur des Wasserstoffs die dimensionslose Zahl
\begin{equation}
\alpha=\frac{e^2}{\hbar c}=\frac{v_1}{c}\approx\frac{1}{137}
\end{equation}
ein~\cite{sommerfeld1916}.
In seinem Gedankenexperiment war $\alpha$ das Verhältnis der Bahngeschwindigkeit auf der innersten Bohrschen Kreisbahn zur Lichtgeschwindigkeit.
Mit der QED wurde daraus eine abstrakte Eichkopplung bei verschwindendem Impulsübertrag, experimentell bestimmt, theoretisch nicht hergeleitet.

Pauli war von der Zahl $137$ besessen.
Er forderte, eine fundamentale Physik dürfe nicht zwanzig freie Parameter anpassen, sondern müsse eine Zahlentheorie des Vakuums sein~\cite{pauli1946}.
Eddington versuchte in den 1930er Jahren, $\alpha^{-1}$ kombinatorisch als $136+1=137$ aus Spinorkomponenten zu gewinnen~\cite{eddington1936}.
Als die Nachkommastellen unmissverständlich bei $137{,}036$ lagen, geriet jede numerische Deduktion in den Geruch der Numerologie.

Das EABC-Framework knüpft an Sommerfeld und Pauli an:
$137$ ist der diskrete Kern eines $E_8$-Torusumlaufs, die Nachkommastellen sind die Hoffbauer-Klammer.
Der Einschluss dieser Klammer ist formal; die Ontologie ``die Natur ist $E_8$'' ist es nicht.

\section{Die algebraische Architektur des $E_8$-Gitters und der Fano-Fluss}

\subsection{Wurzelgeometrie}
Das Fundament ist die größte exzeptionelle Lie-Gruppe $E_8$ mit Rang $8$ und Dimension $248$.
Das Wurzelgitter $\Gamma_8\subset\mathbb{R}^8$ ist das einzige gerade unimodulare Gitter in Dimension~$8$ und maximiert die Packungsdichte~\cite{viazovska2017,conwaysloane}.
Das Wurzelsystem $\Phi(E_8)$ besteht aus $240$ Vektoren der Norm $|x|^2=2$:
\begin{equation}
\Phi(E_8)
=\bigl\{ x\in\mathbb{Z}^8\cup(\mathbb{Z}+\tfrac12)^8
:\textstyle\sum x_i\equiv 0\pmod 2,\; \sum x_i^2=2\bigr\}.
\end{equation}
Sie zerfallen in $112$ Vektoren vom Typ $(\pm 1,\pm 1,0,\dots)$ (das $D_8$-System) und $128$ Halbzahlenvektoren mit gerader Anzahl von Minuszeichen.
Die Theta-Reihe ist die Eisensteinreihe vom Gewicht~$4$,
\begin{equation}
\Theta_{E_8}(\tau)=E_4(\tau)=1+240\sum_{n\ge 1}\sigma_3(n)q^n,\qquad q=e^{2\pi i\tau}.
\end{equation}

\subsection{Oktonionen und Nichtassoziativität}
Der 8-dimensionale Raum ist isomorph zur Divisionsalgebra $\mathbb{O}$ mit Basis $\{1,e_1,\dots,e_7\}$ und Fano-Multiplikation
$e_ie_j=-\delta_{ij}+\sum_k\epsilon_{ijk}e_k$.
Die Oktonionen sind alternativ und strikt nichtassoziativ:
$[x,y,z]=(xy)z-x(yz)\neq 0$ für geeignete Tripel.
Die Automorphismengruppe ist $G_2=\mathrm{Aut}(\mathbb{O})\subset SO(7)$.

\subsection{Die Schale $S^7$ und der Fano-Fluss}
\begin{equation}
S^7:=\{x\in\mathbb{O}:|x|=1\}.
\end{equation}
Die Hopf-Faserung $S^3\hookrightarrow S^7\to S^4$ hat $S^3\cong SU(2)\cong\mathrm{Sp}(1)$ als Einheitsquaternionen.
Im Modell bewegen sich Zustände als gekoppelte Zirkulationen auf Geodäten dieser Schale, geführt durch $\epsilon_{ijk}$.
Das ist die kinematische Bühne der Hoffbauer-Klammer, nicht bereits eine QFT auf $S^7$.

\section{Die Hoffbauer-Klammer}

Der ungestörte Basiswert für den geschlossenen Phasenzyklus auf dem $E_8$-Gitter ist $\alpha_0^{-1}=137$.
Die primäre Projektion der acht Dimensionen auf die Achse von $\mathbb{Q}(\sqrt{5})$ erzeugt den radialen Massendefekt
\begin{equation}
D_r:=\frac{\mathrm{rk}(E_8)}{137\varphi}=\frac{8}{137\varphi}=\frac{4(\sqrt{5}-1)}{137}.
\end{equation}
Mit $\varphi\approx 1{,}618033988749895$ folgt $D_r\approx 0{,}036089575985395$.
Gegen die empirische Fenstermitte $\varepsilon\approx 0{,}03599954$ bleibt
$\Delta_{\mathrm{res}}=D_r-\varepsilon\approx +9{,}00\times 10^{-5}>0$.
(Der Vergleich mit dem CODATA-2022-Zentralwert $0{,}035999177$ ändert das Vorzeichen des Residuums nicht und vergrößert es sogar.)

Hier liegt das Vorzeichen-Paradoxon jeder euklidischen Modellierung.

\begin{proposition}[Ausschluss euklidischer Kopplung]
\label{prop:eukl}
Für $D_t>0$ gilt $\sqrt{D_r^2+D_t^2}\ge D_r>\varepsilon$.
Jeder euklidisch additive Torusanteil treibt $\alpha^{-1}$ aus dem Fenster.
\end{proposition}

\begin{definition}[Hoffbauer-Klammer]
\label{def:hoff}
\[
[D_r,D_t]_{\mathcal{H}}:=\sqrt{D_r^2-D_t^2},\qquad
\alpha^{-1}_{\mathrm{EABC}}:=137+[D_r,D_t]_{\mathcal{H}}.
\]
\end{definition}
Das Minuszeichen ist Vakuumpolarisation: transversale $T^8$-Zirkulation schirmt radiale Trägheit ab.
Lineares Nenner-Feedback $137\mapsto\alpha^{-1}$ ist zu grob und wird nicht benutzt, sobald $c$ feiner als $3$ ist.

\section{Dreistufige Hecke-Filtrierung}

$D_t=1/(c\cdot 137)$ mit $c\in\mathbb{Q}$ aus drei kanonischen Stufen.

\subsection{Stufe 1: Dedekind}
$\mathrm{Spin}(8)$ ist die einzige einfach zusammenhängende Lie-Gruppe mit $\mathrm{Out}(\mathrm{Spin}(8))\cong\mathfrak{S}_3$, Ordnung $6$.
Die zyklische Untergruppe $\mathbb{Z}_3$ permutiert $\mathbf{8}_v\to\mathbf{8}_s\to\mathbf{8}_c\to\mathbf{8}_v$.
Die Diskriminante $\Delta(\tau)=\eta(\tau)^{24}$ hat Gewicht $12$; die Stufe $24$ ist die Wurzelzahl von $D_4\subset E_8$.
Die Dämpfung der Trialität liefert
$c_1=3(1-1/24)=23/8$.
Dann $D_{t,1}=8/(23\cdot 137)=8/3151$ und
$\alpha^{-1}_1\approx 137{,}03600016112$, Residuum $+6{,}21\times 10^{-7}$ gegen die Fenstermitte.

\subsection{Stufe 2: Modulo-12}
Der Cartan-Rang $8$ koppelt an die Periode $12$ der Spitzen:
$N=8\times 12=96$, $c_2=23/8-1/96=275/96$.
(Hier ist $275=5^2\cdot 11$ eine ganzzahlige Beobachtung, kein zusätzliches Axiom.)
$D_{t,2}=96/37675$, $\alpha^{-1}_2\approx 137{,}03599950883$, Residuum $-3{,}12\times 10^{-8}$.

\subsection{Stufe 3: Hecke $T_3$}
Der Raum $S_{12}(\mathrm{SL}_2(\mathbb{Z}))$ wird von $\Delta(\tau)=\sum\tau(n)q^n$ erzeugt, $\tau(3)=252$.
Die Trialitäts-Involution halbiert die Weyl-Kammer: $\kappa=\tau(3)/2=126$.
Die Verschiebung $\delta c=1/(8\cdot 252)=1/2016$ setzt sich mit Stufe~2 zusammen:
\begin{equation}
c_3=\frac{275}{96}+\frac{1}{2016}=\frac{361}{126}=\frac{19^2}{\tau(3)/2}.
\end{equation}
$19$ ist Heegner-Zahl; das ist Kommentar, nicht Fit.
\begin{equation}
D_t=\frac{126}{361\cdot 137}=\frac{126}{49457}.
\end{equation}

\section{Analytischer Wert und Zertifizierung}

\begin{theorem}[Analytische Form]
\label{thm:alpha}
\[
\alpha^{-1}_{\mathrm{EABC}}
=137+\sqrt{\Bigl(\frac{8}{137\varphi}\Bigr)^2-\Bigl(\frac{126}{49457}\Bigr)^2}.
\]
Äquivalent ist der Radikand
\[
\frac{32(3-\sqrt{5})}{137^2}-\frac{15876}{2445994849}.
\]
Numerisch $\alpha^{-1}_{\mathrm{EABC}}\approx 137{,}03599954005601$.
\end{theorem}

\begin{table}[ht]
\centering
\small
\begin{tabular}{@{}lccc@{}}
\toprule
Stufe & $c$ & $\alpha^{-1}$ & Residuum vs.\ $137{,}03599954$ \\
\midrule
$D_r$ allein & $\infty$ & $137{,}036089576$ & $+9{,}04\times 10^{-5}$ \\
Dedekind & $23/8$ & $137{,}036000161$ & $+6{,}21\times 10^{-7}$ \\
Modulo-12 & $275/96$ & $137{,}035999509$ & $-3{,}12\times 10^{-8}$ \\
Hecke $T_3$ & $361/126$ & $137{,}035999540056$ & $+5{,}60\times 10^{-11}$ \\
\midrule
CODATA 2022 & -- & $137{,}035999177(21)$ & $+3{,}63\times 10^{-7}$ \\
\bottomrule
\end{tabular}
\caption{Konvergenz gegen die Fenstermitte; letzte Zeile ist der CODATA-2022-Abstand, kein ppt-Treffer.}
\end{table}

Das $0{,}41\,\mathrm{ppt}$-Residuum gilt gegen $137{,}03599954$, nicht gegen CODATA~2022.

C++ (\texttt{eabc\_alpha.hpp}) benutzt \texttt{save\_state<rounded\_transc\_std<double>>};
das Maschinenintervall hat Breite eines ULP um $137{,}035999540056$.
Lean~4 (Mathlib 4.29) beweist
\begin{equation}
137{,}03599954005<\alpha^{-1}_{\mathrm{EABC}}<137{,}03599954007
\end{equation}
über zwölfstellige $\sqrt{5}$-Zertifikate und \texttt{norm\_num} auf $\mathbb{Q}$~\cite{demoura2021,mathlib}.
Die Identitäten $D_r=4(\sqrt{5}-1)/137$, $D_r^2=32(3-\sqrt{5})/137^2$ und $D_t^2=15876/2445994849$ sind formalisiert.

\section{Planck-Skalen und $G=4$}

Im geometrischen Maß $e=c=1$ ist $\hbar=\alpha^{-1}_{\mathrm{EABC}}$.
Das Postulat $G:=4=\dim_{\mathbb{R}}\mathbb{H}$ ist holografisch motiviert durch $S_{\mathrm{BH}}=A/(4\ell_P^2)$~\cite{bekenstein1973,hawking1975}, nicht der SI-Wert von $G$.
Dann
\begin{align}
\ell_P&=2\sqrt{\alpha^{-1}}\approx 23{,}41247526769,\\
m_P&=\tfrac12\sqrt{\alpha^{-1}}\approx 5{,}85311881692,\\
\ell_P m_P&=\alpha^{-1},\qquad \ell_P/m_P=4,\\
\mathcal{A}_0&=16\alpha^{-1}\approx 2192{,}57599264090,\\
h&=2\pi\hbar\approx 861{,}02257886475.
\end{align}
Es gilt $h_0=2\pi\cdot 137\approx 860{,}79638708<861=\Delta_{41}<h$.
Der Sprung über die 41.\ Dreieckszahl folgt bereits aus $D_r>861/(2\pi)-137\approx 0{,}0324$;
Hecke setzt nur den Rest $h-861\approx 0{,}02257886$.
Euler-Polynom und Heegner-Zahl $41$ sind Beobachtungen, keine Eingangsparameter.
Auf der Einheitsmasse ist $F_g=4F_e$.
Die SI-Hierarchie ist damit nicht gelöst.

\section{Navier--Stokes: Regularitätsvorschlag, kein Clay-Beweis}

Das Millennium-Problem verlangt, ob die dreidimensionalen inkompressiblen Navier--Stokes-Gleichungen für glatte, schnell fallende Anfangsdaten globale glatte Lösungen besitzen oder in endlicher Zeit eine Singularität bilden~\cite{fefferman2000,clay2000}.
Das Beale--Kato--Majda-Kriterium verknüpft den Blow-up mit unintegrierbarer Vortizität
$\int_0^{T^*}\|\omega(\cdot,\tau)\|_\infty\,d\tau=\infty$.

Im Modell ist ein solcher Blow-up ein Signal, dass die kontinuierliche PDE unterhalb einer Gitterlänge unvollständig ist.
Ein kontrahierender Wirbelkern mit Radius $R(t)\to 0$ wird an $R_{\min}=\ell_P=2\sqrt{\alpha^{-1}}$ im geometrischen Maß gestoppt gelesen;
die Hoffbauer-Klammer bleibt strikt kleiner als $D_r$, und die Zirkulation ist durch $h=2\pi\alpha^{-1}$ quantisiert.
Das Kelvin-Theorem $\mathrm{d}\Gamma/\mathrm{d}t=0$ setzt stetig verformbare Kurven voraus; unterhalb der Fläche $\mathcal{A}_0=16\alpha^{-1}$ ist diese Voraussetzung im Modell verletzt.

\begin{remark}
Das ist ein Regularitätsmechanismus im geometrischen Maß, kein Existenz- und Glattheitssatz auf $\mathbb{R}^3$.
$\ell_P\approx 23{,}41$ ist dimensionslos, nicht die SI-Planck-Länge.
Machine-learning-gestützte Blow-up-Kandidaten der Euler- oder Navier--Stokes-Literatur werden hier nicht geprüft.
Die Clay-Prämie ist durch diese Note nicht beansprucht.
\end{remark}

\section{Yang--Mills: Modell-Massenlücke, kein Clay-Beweis}

Das Clay-Problem verlangt die Konstruktion einer nichtabelschen Eichentheorie auf $\mathbb{R}^4$ im Sinne der Wightman- oder Osterwalder--Schrader-Axiome und den Beweis einer strikt positiven Lücke zwischen Vakuum und erstem Glueball~\cite{jaffewitten2000}.
Ein Kontinuums-Pfadintegral $\int\mathcal{D}A\,e^{-S}$ wird im EABC-Modell nicht konstruiert.
Das Vakuum wird als Zustandssumme $\Theta_{E_8}=E_4$ gelesen -- ebenfalls eine Analogie, keine OS-Rekonstruktion.

Definiert man $\Delta_{\mathrm{YM}}:=m_P[D_r,D_t]_{\mathcal{H}}$, so ist $\Delta_{\mathrm{YM}}>0$ äquivalent zu $D_r\neq D_t$.
Gleichheit erzwänge $\varphi=8\cdot 49457/(137\cdot 126)\in\mathbb{Q}$, Widerspruch zur Irrationalität von $\sqrt{5}$.

\begin{remark}
Die Irrationalität von $\varphi$ schützt den Radikanden der Hoffbauer-Klammer, nicht die Massenlücke der Kontinuums-Yang--Mills-Theorie.
Ein Vergleich $M(0^{++})\approx 1{,}7\,\mathrm{GeV}$ mit Gitter-QCD ist hier nicht hergeleitet.
\end{remark}

\section{Kosmologie: Probeformeln}

Die folgenden Ausdrücke sind spekulativ.
Ehrliche Auswertung, ohne Anpassung an Planck-2018-Zentralwerte:
\begin{align}
\Omega_\Lambda^{\mathrm{(Probe)}}&=1-\varphi^{-2}-1/(24\varphi)\approx 0{,}5923
\quad\text{(Planck 2018: }0{,}6847\pm 0{,}0073\text{)},\\
n_s^{\mathrm{(Probe)}}&=1-1/24-1/(24\varphi^3)\approx 0{,}9485
\quad\text{(Planck 2018: }0{,}9649\pm 0{,}0042\text{)}.
\end{align}
Die in Entwürfen behauptete Identität mit $0{,}6847$ bzw.\ $0{,}9649$ ist arithmetisch falsch.

Der Hubble-Faktor
\begin{equation}
\Bigl(\frac{12}{\sqrt{\alpha^{-1}}}\Bigr)^3\approx 1{,}07719
\end{equation}
liefert $67{,}40\times 1{,}07719\approx 72{,}60\,\mathrm{km\,s^{-1}\,Mpc^{-1}}$.
Das ist eine Skalenprobe, keine Lösung der Hubble-Spannung:
CMB und SH0ES messen verschiedene Observable, nicht $S^7$ versus Chandrasekhar-Volumen als bewiesenen Dimensionssprung~\cite{planck2018,riess2022}.

\section{Familien, Higgs, Spektroskopie}

$\mathrm{Out}(\mathrm{Spin}(8))\cong\mathfrak{S}_3$ permutiert $\mathbf{8}_v,\mathbf{8}_s,\mathbf{8}_c$.
Drei Familien als $\mathbb{Z}_3$-Orbit und Higgs als Radialmode von $S^7$ mit $\mathbb{H}^\perp\cong\mathbb{R}^4$ sind geometrische Lesarten~\cite{adams}.
\begin{equation}
\frac{1}{\sqrt{2\varphi}}\approx 0{,}55589,\qquad
v\cdot\frac{1}{\sqrt{2\varphi}}\approx 136{,}87\,\mathrm{GeV}\neq 125{,}25\,\mathrm{GeV}.
\end{equation}
Die Kette $246{,}22\times 0{,}5559\approx 125{,}3$ ist falsch.
$\varphi^{-40}$ reproduziert $v\approx 246\,\mathrm{GeV}$ nicht.

Der Schwinger-Term $a^{(1)}=\alpha/(2\pi)=1/(2\pi\alpha^{-1})$ ist in den Modell-Einheiten $1/h$.
Die Subtraktion $\Delta_{\mathrm{EABC}}/(2\pi(\alpha^{-1})^2)$ trifft den Gabrielse-Wert $0{,}00115965218073$ nicht~\cite{gabrielse2023}.
Eine geschlossene Lösung von $a_e$ oder $a_\mu$ wird hier nicht behauptet.

\section{Gesicherte versus offene Größen}

\begin{table}[ht]
\centering
\small
\begin{tabular}{@{}llll@{}}
\toprule
Größe & Status & Ausdruck & Wert \\
\midrule
$\alpha^{-1}$ & Lean/C++ & $137+\sqrt{D_r^2-D_t^2}$ & $137{,}03599954005601$ \\
$D_r$ & Lean & $4(\sqrt{5}-1)/137$ & $0{,}036089575985395$ \\
$D_t$ & Lean & $126/49457$ & $0{,}002547667670906$ \\
$\hbar$ & Postulat $e=c=1$ & $\alpha^{-1}$ & $137{,}03599954005601$ \\
$h$ & $2\pi\hbar$ & $2\pi\alpha^{-1}$ & $861{,}02257886475$ \\
$G$ & Postulat & $4$ & $4$ \\
$\ell_P,m_P$ & aus $G=4$ & $2\sqrt{\alpha^{-1}}$, $\tfrac12\sqrt{\alpha^{-1}}$ & $23{,}41247527$, $5{,}85311882$ \\
\midrule
NS/YM Clay & offen & Regularitätslesart & -- \\
$\Omega_\Lambda,n_s,m_H,a_e$ & nicht geschlossen & Probeformeln weichen ab & -- \\
CODATA-2022-$1\sigma$ & nicht getroffen & Abstand $3{,}63\times 10^{-7}$ & $\sim 17\sigma$ \\
\bottomrule
\end{tabular}
\caption{Was zertifiziert ist, und was Lesart bleibt.}
\end{table}

\section{Epilog}

Die Hoffbauer-Klammer
\begin{equation}
[D_r,D_t]_{\mathcal{H}}
=\sqrt{\Bigl(\frac{8}{137\varphi}\Bigr)^2-\Bigl(\frac{126}{49457}\Bigr)^2}
\end{equation}
ist der gesicherte arithmetische Kern:
sie verbindet $\mathbb{Q}(\sqrt{5})$ mit $\tau(3)$ und schließt $\alpha^{-1}$ in ein Sub-PPB-Fenster.
Das Kontinuum bleibt für makroskopische Skalen die richtige effektive Beschreibung.
Wo Theorien divergieren, ist eine Gittergrenze \emph{denkbar} -- sie ist durch die vorliegende Note nicht als Lösung der Millennium-Probleme oder als Elimination des Standardmodells bewiesen.

\begin{thebibliography}{99}

\bibitem{newton1687}
I.~Newton,
\emph{Philosophi{\ae} Naturalis Principia Mathematica},
London, 1687.

\bibitem{planck1900}
M.~Planck,
\emph{Zur Theorie des Gesetzes der Energieverteilung im Normalspektrum},
Verh.\ Dtsch.\ Phys.\ Ges.\ \textbf{2} (1900), 237--245.

\bibitem{sommerfeld1916}
A.~Sommerfeld,
\emph{Zur Quantentheorie der Spektrallinien},
Ann.\ Phys.\ \textbf{51} (1916), 1--94.

\bibitem{eddington1936}
A.~S.~Eddington,
\emph{Relativity Theory of Protons and Electrons},
Cambridge University Press, 1936.

\bibitem{pauli1946}
W.~Pauli,
\emph{Exclusion Principle and Quantum Mechanics},
Nobel Lecture, 1946.

\bibitem{dirac1975}
P.~A.~M.~Dirac,
in: \emph{Directions in Physics}, Wiley, 1978.

\bibitem{penrose1965}
R.~Penrose,
Gravitational collapse and space-time singularities,
Phys.\ Rev.\ Lett.\ \textbf{14} (1965), 57--59.

\bibitem{hawking1970}
S.~W.~Hawking, R.~Penrose,
The singularities of gravitational collapse and cosmology,
Proc.\ R.\ Soc.\ Lond.\ A \textbf{314} (1970), 529--548.

\bibitem{bekenstein1973}
J.~D.~Bekenstein,
Black holes and entropy,
Phys.\ Rev.\ D \textbf{7} (1973), 2333--2346.

\bibitem{hawking1975}
S.~W.~Hawking,
Particle creation by black holes,
Commun.\ Math.\ Phys.\ \textbf{43} (1975), 199--220.

\bibitem{codata2022}
P.~J.~Mohr, D.~B.~Newell, B.~N.~Taylor, E.~Tiesinga,
CODATA recommended values of the fundamental physical constants: 2022,
Rev.\ Mod.\ Phys.\ \textbf{97} (2025), 025002.

\bibitem{conwaysloane}
J.~H.~Conway, N.~J.~A.~Sloane,
\emph{Sphere Packings, Lattices and Groups},
Springer, 3.\ Aufl., 1999.

\bibitem{viazovska2017}
M.~S.~Viazovska,
The sphere packing problem in dimension 8,
Ann.\ of Math.\ \textbf{185} (2017), 991--1015.

\bibitem{adams}
J.~F.~Adams,
\emph{Lectures on Exceptional Lie Groups},
University of Chicago Press, 1996.

\bibitem{clay2000}
J.~Carlson, A.~Jaffe, A.~Wiles (Hrsg.),
\emph{The Millennium Prize Problems},
CMI/AMS, 2006.

\bibitem{fefferman2000}
C.~L.~Fefferman,
Existence and smoothness of the Navier--Stokes equation,
in~\cite{clay2000}.

\bibitem{jaffewitten2000}
A.~Jaffe, E.~Witten,
Quantum Yang--Mills theory,
in~\cite{clay2000}.

\bibitem{planck2018}
Planck Collaboration,
Planck 2018 results.\ VI.\ Cosmological parameters,
Astron.\ Astrophys.\ \textbf{641} (2020), A6.

\bibitem{riess2022}
A.~G.~Riess et al.,
A comprehensive measurement of the local value of the Hubble constant,
Astrophys.\ J.\ Lett.\ \textbf{934} (2022), L7.

\bibitem{gabrielse2023}
X.~Fan, T.~G.~Myers, B.~A.~D.~Sukra, G.~Gabrielse,
Measurement of the electron magnetic moment,
Phys.\ Rev.\ Lett.\ \textbf{130} (2023), 071801.

\bibitem{demoura2021}
L.~de~Moura, S.~Ullrich,
The Lean~4 Theorem Prover,
in: CADE 28, LNCS 12699, Springer, 2021.

\bibitem{mathlib}
The mathlib Community,
The Lean mathematical library,
in: CPP 2020, ACM, 2020.

\end{thebibliography}

\end{document}
